\section{Background and motivation}
\label{sec:motivation}
\paragraph{Transfer a debugging decision across repositories.}
The Django case (Trac \#33018) mishandles empty query results under negation.
A failed \texttt{python-statemachine} attempt had spent its effort setting
up a full Django application. Its lesson was to isolate the bug with a
minimal reproduction before investing in framework setup. A successful
\texttt{sievelib} trajectory also supported reproducing the failure before
editing. Together, these experiences suggest an investigation strategy and
warn against an expensive detour.

\paragraph{Apply the lesson to the current code.}
The history-informed agent followed the exception through query compilation
to \texttt{Exists.as\_sql()}, corrected its handling, and added a regression
test. The no-memory attempt ended with broad formatting changes and no core
fix. Memory supplied the strategy; the agent worked out the local implementation.

Figure~\ref{fig:concept} illustrates how retrieved requirements guide repair
checks in our behavioral studies.

\begin{figure}[!htbp]
\centering
\includegraphics[width=\linewidth]{figures/contextgraph-design/fig1-concept.pdf}
\caption{Behavioral study of a retrieved requirement in a Django repair. (a) The agent receives
a requirement as text. (b) Linked context supports a check that combines the
historical range type $R$ with the current query. Execution feedback guides
the next edit.}
\label{fig:concept}
\end{figure}

\section{ContextGraph: memory for planning and error recovery}
\label{sec:method}
ContextGraph connects three designs: experience abstraction makes past
decisions transferable, graph retrieval finds related lessons, and memory
access brings them into planning and error recovery
(Figure~\ref{fig:experience-memory}).

\begin{figure}[!htbp]
\centering
\includegraphics[width=\linewidth]{figures/figure3-design-v1/fig3-contextgraph.pdf}
\caption{ContextGraph's experience-memory architecture. (a) Past trajectories
supply transferable strategies and warnings linked to source context.
(b) The agent retrieves guidance before planning and after errors, then applies
it through investigation, edits, and tests. The graph shows illustrative links.}
\label{fig:experience-memory}
\end{figure}

\subsection{Experience abstraction: retain the useful decision}
\label{sec:experience-abstraction}
From task descriptions, actions, errors, and outcomes, a language model
extracts lessons about debugging, testing, navigation, and failed approaches.
Abstraction removes source-specific names while retaining why the action
helps: the failed framework setup becomes advice to isolate the bug first.

The graph links strategies to source experiences and recurring error patterns.
Similar strategies share a rule, bringing together observations across
repositories. Strategies describe productive actions; warnings identify
unproductive approaches. Links to source outcomes help the agent interpret both.

\subsection{Graph retrieval: find relevant strategies and warnings}
\label{sec:experience-retrieval}
A query describes the task and the agent's situation. Semantic search finds
related lessons; keyword search matches error terms. A graph walk follows
error patterns through shared rules and source strategies to past trajectories
(Figure~\ref{fig:experience-memory}). Different errors can reach the same rule.
The retriever merges semantic, keyword, and graph candidates, selecting a
small, diverse set of lessons.

Returned memory contains actions, warnings, and source context. The agent
interprets this guidance against the current code and develops the local
reproduction and patch. Appendix~\ref{sec:agentic-retrieval} describes
iterative agent search over the experience pool.

\subsection{Memory access: before planning and after errors}
\label{sec:experience-access}
\paragraph{Before generating a plan.}
The agent queries memory with the issue and available code context.
Strategies suggest where to investigate, how to reproduce the problem, and
what to check first; warnings identify costly approaches to avoid.

\paragraph{After encountering an error.}
A failed command, reproduction, or test supplies a more specific query:
the error, recent action, relevant files, and current hypothesis. The same
graph now retrieves guidance on alternative causes and next actions.

\paragraph{From guidance to action.}
The agent applies advice through repository tools: narrowing a search,
tracing an exception, or adding a test. New observations guide further
retrieval. Completed trajectories can enter the same abstraction pipeline,
extending the experience available to later tasks. Figure~\ref{fig:behavior-workflow}
shows how the behavioral studies add task-specific checks to examine the
application of retrieved requirements (Section~\ref{sec:evidence}).

\begin{figure}[!htbp]
\centering
\includegraphics[width=\linewidth]{figures/contextgraph-design/fig2-overview.pdf}
\caption{Workflow for the behavioral studies. (A) Historical
repairs supply requirements and checked examples linked in memory.
(B) Retrieval selects relevant sources; graph relations support composition
with the current operation. Execution feedback guides edits before verification.}
\label{fig:behavior-workflow}
\end{figure}
